\section{Transformers}
\label{sec:transformers}

A convolution uses the same filter at every image location, but a single filter
only sees a local neighborhood. In a sentence, the useful context for a word may
be many positions away: in ``the food was delicious, but the service was slow,''
the two adjectives refer to different nouns. A transformer lets each position
combine information from every position. The strength of each connection is
computed from the current input rather than fixed in advance. Its learned
parameters are shared across positions, so the same layer can process sequences
of different lengths.

\subsection{Tokens and Embeddings}

Before a network can process text, a tokenizer divides it into discrete units.
Using complete words alone makes the vocabulary enormous and leaves new names or
spellings without a representation. Subword tokenization instead builds a
vocabulary of frequent character sequences. Repeatedly merging common adjacent
pieces can turn characters into a mixture of fragments and whole words. Related
forms such as ``walk,'' ``walked,'' and ``walking'' can then share pieces, while
an unfamiliar word can still be assembled from smaller units.

Give each of the $|\mathcal V|$ vocabulary entries an integer index. A learned
embedding table $E\in\mathbb R^{|\mathcal V|\times d}$ maps token $t_i$ to
$x_i=E_{t_i,:}\in\mathbb R^d$. The same table is used at every position. This
vector is a starting representation, not a meaning fixed for all contexts: the
token ``bank'' has the same input vector in ``river bank'' and ``bank account.''
Attention can give those occurrences different outputs by combining them with
their neighbors.

Token order must also be supplied. Without position information, permuting the
input tokens merely permutes the outputs of self-attention: it cannot tell ``the
man ate the fish'' from ``the fish ate the man'' at a given token position.
An \emph{absolute positional encoding} adds a vector $p_i$ to each embedding,
giving the layer input $\tilde x_i=x_i+p_i$. One fixed choice uses different
sinusoidal frequencies in successive coordinate pairs:
\begin{equation}
  p_{i,2j}=\sin\!\left(\frac{i}{10000^{2j/d}}\right),\qquad
  p_{i,2j+1}=\cos\!\left(\frac{i}{10000^{2j/d}}\right).
  \label{eq:transformer-position}
\end{equation}
\noindent
Positions may instead have learned vectors. A relative position scheme can add
a learned bias $b_{i-j}$ to the attention score from position $i$ to position
$j$. Both approaches make order available; the first labels individual
positions, while the second labels their separation.

\subsection{Self-Attention}

Suppose $X\in\mathbb R^{N\times d}$ holds $N$ token vectors as rows, including
their position information. Each token plays three learned roles. Its
\emph{query} describes what information it seeks, its \emph{key} describes how
it can be found, and its \emph{value} carries the information to contribute.
Shared linear projections produce all three:
\begin{equation}
  Q=XW^Q,\qquad K=XW^K,\qquad V=XW^V,
  \label{eq:transformer-qkv}
\end{equation}
\noindent
where $W^Q,W^K\in\mathbb R^{d\times d_k}$ and
$W^V\in\mathbb R^{d\times d_v}$. Query $q_i$ compares with every key $k_j$.
The resulting scores pass through a softmax over $j$:
\begin{equation}
  s_{ij}=\frac{q_i k_j^{\mathsf T}}{\sqrt{d_k}},\qquad
  a_{ij}=\frac{\exp(s_{ij})}{\sum_{\ell=1}^{N}\exp(s_{i\ell})},\qquad
  y_i=\sum_{j=1}^{N}a_{ij}v_j.
  \label{eq:transformer-attention}
\end{equation}
\noindent
The nonnegative weights in each row sum to one. Thus $y_i$ is a weighted
combination of values, with weights chosen separately for query $i$. The
$1/\sqrt{d_k}$ factor keeps dot products from growing with key dimension when
their coordinates have roughly unit variance; otherwise softmax can become so
peaked that its gradients are small.

For a small calculation, take $d_k=2$, $q=(1,0)$, keys $k_1=(1,0)$ and
$k_2=(0,1)$, and values $v_1=(2,0)$ and $v_2=(0,2)$. The scores are
$(1/\sqrt{2},0)$, so the softmax weights are approximately $(0.670,0.330)$.
The output is $0.670(2,0)+0.330(0,2)=(1.340,0.660)$. Changing the query
changes this mixture even when the two values stay the same.

All queries can be processed together. With softmax applied separately to each
row, Equation~\ref{eq:transformer-attention} becomes
\begin{equation}
  A=\operatorname{softmax}_{\mathrm{row}}\!\left(
      \frac{QK^{\mathsf T}}{\sqrt{d_k}}\right),\qquad
  Y=AV.
  \label{eq:transformer-matrix-attention}
\end{equation}
\noindent
Here $A\in\mathbb R^{N\times N}$ and $Y\in\mathbb R^{N\times d_v}$.
Every output can use every input after one layer. Forming all pairwise scores
takes $O(N^2d_k)$ arithmetic and $O(N^2)$ attention storage, which becomes
costly for long sequences. Unlike a recurrent network, the score calculations
for different positions can run in parallel.

One attention pattern need not serve every relationship. \emph{Multihead
attention} learns separate projections for $h$ heads, computes
$H_r=\operatorname{Attention}(XW_r^Q,XW_r^K,XW_r^V)$ for each head $r$, and
concatenates their outputs:
\begin{equation}
  \operatorname{MHA}(X)=
  \operatorname{Concat}(H_1,\ldots,H_h)W^O.
  \label{eq:transformer-mha}
\end{equation}
\noindent
Each head can weight a different set of tokens. If its value width is $d_v$,
then $W^O\in\mathbb R^{hd_v\times d}$ returns to the model width $d$.
Multiple heads do not change the quadratic dependence on sequence length.

\subsection{Transformer Layers}

A transformer layer combines multihead attention with a small network applied
independently to each token. Both parts have residual connections, as in
Section~\ref{sec:resnet}, and layer normalization. The attention part mixes
information \emph{between} tokens; the feed forward part transforms features
\emph{within} each token. Figure~\ref{fig:transformer-layer} shows one layer.

\begin{figure}[ht]
  \centering
  \begin{tikzpicture}[
    box/.style={draw, rounded corners=2pt, minimum width=2.25cm,
      minimum height=0.85cm, align=center, font=\small},
    link/.style={-{Latex[length=2mm]}, thick}
  ]
    \node[box] (input) at (0,0) {Input $X$};
    \node[box] (attn) at (2.7,0) {Multihead\\attention};
    \node[box] (norm1) at (5.4,0) {Add \& norm};
    \node[box] (ffn) at (8.1,0) {Feed forward};
    \node[box] (norm2) at (10.8,0) {Add \& norm};
    \node[box] (output) at (13.5,0) {Output $Z$};
    \draw[link] (input) -- (attn);
    \draw[link] (attn) -- (norm1);
    \draw[link] (norm1) -- (ffn);
    \draw[link] (ffn) -- (norm2);
    \draw[link] (norm2) -- (output);
    \draw[link] (input.north) -- ++(0,1.05) -| (norm1.north);
    \draw[link] (norm1.south) -- ++(0,-1.05) -| (norm2.south);
  \end{tikzpicture}
  \caption{One transformer layer. The upper and lower paths add each
  sublayer's input to its output before normalization. The feed forward
  network acts on each token separately.}
  \label{fig:transformer-layer}
\end{figure}

\noindent
For the layout shown, the operations can be written
\begin{equation}
  H=\operatorname{LN}\bigl(X+\operatorname{MHA}(X)\bigr),\qquad
  Z=\operatorname{LN}\bigl(H+\operatorname{FFN}(H)\bigr).
  \label{eq:transformer-layer}
\end{equation}
\noindent
Layer normalization computes a mean and variance across the $d$ features of
\emph{each token}, then applies learned scale and shift parameters. It does not
require batch statistics. The same two dense layers and biases are reused at
every position; for token row $h_i$ an example is
\begin{equation}
  \operatorname{FFN}(h_i)=
  \operatorname{ReLU}(h_iW_1+b_1)W_2+b_2,
  \quad W_1\in\mathbb R^{d\times d_{\mathrm{ff}}},\quad
  W_2\in\mathbb R^{d_{\mathrm{ff}}\times d}.
  \label{eq:transformer-ffn}
\end{equation}
\noindent
Its input and output both have width $d$, so the second residual addition is
valid. Stacking layers lets each token repeatedly gather context and transform
what it gathered. A transformer \emph{encoder} uses such layers to represent
an entire input sequence. A \emph{decoder} uses transformer layers to produce
an output sequence, restricting attention when it must predict tokens in order.

\subsection{Encoders and Transfer Learning}

An encoder reads all input tokens at once. Its attention is bidirectional: an
output token may use words both before and after its position. This makes an
encoder useful when the whole input is available and the goal is to understand
it. BERT is an encoder example. During masked-language-model pretraining, some
input tokens are hidden and the model predicts their identities from the
surrounding tokens. A classifier applied to the output at a masked position
produces a probability for each vocabulary entry; training increases the
probability of the original entry. The target is a \emph{token identity}, not
the hidden token's embedding vector. Because the same layers must recover many
different missing tokens, their outputs learn contextual representations.

These representations can support several tasks. For sentence classification,
a special $[\mathrm{CLS}]$ token is inserted and its final output is passed to
a small classification head. Sentiment classification is one example. A pair
of sentences can be encoded together for a relation decision. For token
classification, a head reads each token's final output, as in named-entity
recognition, where words are labeled as person, place, organization, or none.
Span annotation marks a contiguous range of tokens, for example by predicting
its start and end. The output used by the head therefore depends on whether
the target describes a whole input, each token, or a span.

Pretraining first learns reusable features on a large source dataset; a
labeled target dataset then adapts the model. In \emph{fine-tuning}, the new
task's loss updates some or all pretrained weights as well as the task head.
Alternatively, a \emph{linear probe} freezes the encoder and trains only a
linear head. Freezing reduces the number of trained parameters, but does not
remove the encoder at test time. Fine-tuning does not require freezing all
layers except the head. Transfer can help even when source and target tasks
differ because early visual filters or general contextual features are useful
across tasks; the final decision rule may still need to be learned for the new
labels. This is why a pretrained initialization can be better than random
weights without promising that every feature or class decision transfers.

\subsection{Decoders and Cross-Attention}

A decoder generates one token at a time. For tokens $t_1,\ldots,t_N$, an
autoregressive model factors their joint probability as
\begin{equation}
  p(t_1,\ldots,t_N)
  =\prod_{i=1}^{N}p(t_i\mid t_1,\ldots,t_{i-1}).
  \label{eq:transformer-autoregressive}
\end{equation}
\noindent
A GPT-style decoder implements these conditional distributions. The output at
position $i$ predicts the next token from a softmax over the vocabulary.
\emph{Causal masking} prevents a query at position $i$ from attending to
future positions $j>i$: their attention scores are set to $-\infty$ before
softmax. Without that mask, training could read the answer it is meant to
predict. A start token and a shifted target sequence align each output with
its next-token target. Known target tokens can be processed in parallel during
training because the mask hides later tokens. At inference, each new token
must be chosen before the next step, so generation remains sequential.
Providing a few examples in the input prompt can guide a pretrained decoder
without changing its weights; this is in-context or few-shot learning, not
fine-tuning.

For translation, the source sentence is already known, but the target sentence
must be generated in order. An encoder first forms source representations
$E\in\mathbb R^{N_e\times d}$. The decoder uses masked self-attention over its
target prefix and then \emph{cross-attention}: its target states
$D\in\mathbb R^{N_d\times d}$ supply queries, while the encoder states supply
keys and values,
\begin{equation}
  Q=DW^Q,\quad K=EW^K,\quad V=EW^V,\qquad
  C=\operatorname{softmax}_{\mathrm{row}}\!\left(
       \frac{QK^{\mathsf T}}{\sqrt{d_k}}\right)V.
  \label{eq:transformer-cross-attention}
\end{equation}
\noindent
The cross-attention weights have shape $N_d\times N_e$: each target position
can select information from every source position. The causal mask applies to
target self-attention, not to these known source states. Translation training
again predicts each next target token; at inference, the decoder uses its
previously generated target tokens. Table~\ref{tab:transformer-roles} separates
the three common layouts.

\begin{table}[ht]
  \centering
  \begin{tabular}{@{}p{3.0cm}p{5.5cm}p{4.9cm}@{}}
    \toprule
    Layout & Information available & Typical output \\
    \midrule
    Encoder & all input tokens & sequence or token labels \\
    Decoder & past target tokens only & next-token probabilities \\
    Encoder--decoder & source and past target tokens & translated sequence \\
    \bottomrule
  \end{tabular}
  \caption{Which tokens each transformer layout can read when forming an
  output. The attention mask determines whether future target tokens are
  hidden.}
  \label{tab:transformer-roles}
\end{table}

\noindent
Full attention over $N$ tokens uses the $N\times N$ weight matrix from
Equation~\ref{eq:transformer-matrix-attention}. Its quadratic storage and
compute become expensive for long text or dense images. Restricting a token to
a local neighborhood, using sparse connections, or introducing selected global
tokens reduces the number of interactions, but also changes which pairs can
exchange information directly. A decoder's triangular causal mask prevents
future leakage; by itself it does not remove the quadratic cost of full
attention over the allowed prefix.

\subsection{Vision and Multiscale Transformers}

An image has many more pixels than a typical sentence has words. Treating each
pixel as a token makes global attention expensive: a $224\times224$ image has
$50{,}176$ pixel positions, so its attention matrix has over $2.5$ billion
entries per head. Early pixel-generating transformer decoders such as ImageGPT
therefore operated on small images. A Vision Transformer (ViT) reduces the
sequence length by dividing an $H\times W\times C$ image into nonoverlapping
$P\times P$ patches. Assuming $P$ divides both dimensions, the number of
patches is
\begin{equation}
  N=\frac{H}{P}\frac{W}{P}=\frac{HW}{P^2},\qquad
  f_i\in\mathbb R^{P^2C},\qquad
  z_i=f_iW^E+p_i\in\mathbb R^d,
  \label{eq:vit-patches}
\end{equation}
\noindent
where $f_i$ is a flattened patch, $W^E\in\mathbb R^{P^2C\times d}$ is shared
across patches, and $p_i$ is a learned positional vector. The linear patch
projection is equivalent to a convolution with kernel size and stride $P$,
but the following encoder uses attention rather than local convolutions.
A learned $[\mathrm{CLS}]$ vector is prepended to the $N$ patch tokens, so
the encoder processes $N+1$ tokens. Its final class-token output passes
through a classification head and softmax. Patch positions must be encoded:
without them, rearranging patches would only rearrange their outputs and the
class token could not recover where image content appeared.

For a $224\times224$ RGB image with $P=16$, each flattened patch has
$16^2\cdot3=768$ values and there are $(224/16)^2=196$ patches, or $197$
tokens with the class token. A head's attention matrix is then
$197\times197$, instead of $50{,}176\times50{,}176$ for pixel tokens. This
calculation explains why patch size matters: smaller patches retain finer
detail but increase attention cost quadratically with their count.

Unlike a convolutional network, a plain ViT has no architectural requirement
that nearby pixels interact first or that shifting an image shifts its feature
map in the same way. Any patch can attend to any other patch in the first
layer. Convolutions instead start with local neighborhoods and build a larger
receptive field across layers. This built-in locality and translation
equivariance help CNNs learn from limited images. ViT must learn useful
spatial relationships from data, but large-scale pretraining can make its
global interactions effective. Attention distance makes the contrast visible:
the average distance between a query patch and its attended patches can be
large even in early ViT layers, while a small convolution initially sees only
its local neighborhood. For a new classification task, a pretrained ViT can
be fine-tuned, or its encoder can be frozen and a new head trained.

Plain ViT keeps one patch scale throughout. A shifted-window transformer
(Swin) instead computes attention within small windows of the patch grid.
The next block shifts the window partition so patches previously separated by
a window boundary can exchange information. Periodic \emph{patch merging}
combines each $2\times2$ group of neighboring patch features and projects it
to a new channel width. Spatial resolution falls while the receptive field
grows, giving a hierarchy more like a CNN. If each attention window contains
at most $M$ tokens, its attention storage over all windows is $O(NM)$ rather
than $O(N^2)$ for global attention, when $M$ is fixed. The shift is essential:
without it, repeated window attention cannot pass information across window
boundaries at the same scale.
The self-supervised ViT case DINOv2 is explained in
Appendix~\ref{app:foundation}.
